\chapter{Propositional formulas and the covering partition}\label{sec:coverage}

This chapter develops, from first principles, the small amount of
propositional logic that the certificate rests on, and then the combinatorial
heart of the matter: the $98{,}536$ cells of the distributed certificate
partition the whole space of valuations.  No background in logic is assumed;
every notion is a finite combinatorial object over the two-element set
$\{0,1\}$.  (In Lean, $\{0,1\}$ is the type \texttt{Bool} with elements
\texttt{false}, \texttt{true}; we write $0$ and $1$.)

\section{Propositional formulas from first principles}

\begin{definition}[Variables and literals]\label{def:cnf-literal}
Fix a countable index set $V$ of \emph{variables}; concretely $V = \{0, 1, 2, \dots\}$.  A
\emph{literal} is a pair $(x, \varepsilon) \in V \times \{0,1\}$, written
$x$ when $\varepsilon = 1$ and $\neg x$ when $\varepsilon = 0$: an assertion
that the variable $x$ takes the value $\varepsilon$.
\lean{Std.Sat.Literal}\leanok
\end{definition}

\begin{definition}[Clauses and formulas]\label{def:cnf-formula}
A \emph{clause} is a finite list of literals, read \emph{disjunctively}
(``at least one of these assertions holds'').  A \emph{formula in
conjunctive normal form} (a \emph{CNF formula}) is a finite family
$F = (C_1, \dots, C_m)$ of clauses, read \emph{conjunctively} (``every
clause holds'').
\lean{Std.Sat.CNF, Std.Sat.CNF.Clause}\leanok
\uses{def:cnf-literal}
\end{definition}

\begin{definition}[Valuations and satisfaction]\label{def:cnf-valuation}
A \emph{valuation} is a function $a \colon V \to \{0,1\}$, assigning a value
to every variable.  A valuation $a$ \emph{satisfies}
\begin{itemize}
\item the literal $(x, \varepsilon)$ iff $a(x) = \varepsilon$;
\item a clause $C$ iff it satisfies \emph{at least one} literal of $C$;
\item a CNF formula $F$ iff it satisfies \emph{every} clause of $F$.
\end{itemize}
In Lean these are the evaluation maps sending a valuation and a
clause/formula to $\{0,1\}$; satisfaction means evaluating to $1$.
\lean{Std.Sat.CNF.eval, Std.Sat.CNF.Clause.eval}\leanok
\uses{def:cnf-formula}
\end{definition}

\begin{definition}[Satisfiability]\label{def:cnf-unsat}
A CNF formula $F$ is \emph{satisfiable} if some valuation satisfies it, and
\emph{unsatisfiable} if no valuation does.  Unsatisfiability of a suitable
formula is the finitary statement to which the non-existence theorem of this
project is reduced.
\lean{Std.Sat.CNF.Sat, Std.Sat.CNF.Unsat}\leanok
\uses{def:cnf-valuation}
\end{definition}

\begin{definition}[Cubes: partial valuations]\label{def:cnf-cube}
A \emph{cube} is a finite list of literals read \emph{conjunctively}
(``all of these assertions hold'') --- dually to a clause.  A valuation
satisfies a cube $c$ iff it satisfies every literal of $c$; thus a cube
whose variables are pairwise distinct is the same thing as a \emph{partial
valuation} with finite domain, and its satisfaction set is the set of all
total valuations extending it.
\lean{ConwayO7.Cube, ConwayO7.Cube.eval}\leanok
\uses{def:cnf-literal, def:cnf-valuation}
\end{definition}

\begin{definition}[Adjoining a cube to a formula]\label{def:cnf-withcube}
For a CNF formula $F$ and a cube $c$, the formula $F \wedge c$ is obtained
by adjoining to $F$ one singleton clause $\{\ell\}$ per literal $\ell$ of
$c$.  This is the formula whose refutation certifies that $F$ has no
satisfying valuation \emph{inside the cell described by $c$}.
\lean{ConwayO7.withCube}\leanok
\uses{def:cnf-formula, def:cnf-cube}
\end{definition}

\begin{lemma}[Semantics of $F \wedge c$]\label{lem:cnf-withcube-eval}
A valuation satisfies $F \wedge c$ if and only if it satisfies both $F$ and
the cube $c$.
\lean{ConwayO7.eval_withCube}\leanok
\uses{def:cnf-withcube, def:cnf-valuation}
\end{lemma}
\begin{proof}[Proof sketch]
Induction on the list $c$: a singleton clause $\{\ell\}$ is satisfied
exactly when the literal $\ell$ is, and adjoining a clause intersects the
satisfaction sets.
\end{proof}

\begin{theorem}[Composition of refutations]\label{thm:cnf-unsat-of-cubes}
Let $F$ be a CNF formula and let $c_1, \dots, c_N$ be cubes whose
satisfaction sets \emph{cover} the space of valuations (every valuation
satisfies some $c_i$).  If $F \wedge c_i$ is unsatisfiable for every $i$,
then $F$ is unsatisfiable.
\lean{ConwayO7.unsat_of_cubes}\leanok
\uses{def:cnf-unsat, def:cnf-withcube}
\end{theorem}
\begin{proof}[Proof sketch]
\uses{lem:cnf-withcube-eval}
Any valuation satisfying $F$ satisfies some $c_i$, hence by
Lemma~\ref{lem:cnf-withcube-eval} it satisfies $F \wedge c_i$,
contradicting the $i$-th refutation.
\end{proof}

\begin{definition}[The audited text convention]\label{def:cnf-dimacs}
The distributed files render literals as signed integers: the literal
$(v, 1)$ (with $v$ counted from $0$) is written as the integer $v+1$, and
$(v, 0)$ as $-(v+1)$.  A clause is a line of literals terminated by $0$; a
cube is a line of literals with no terminator.  The parsing maps from text
to formulas and cubes, and the serialisation maps in the other direction,
are fixed once in a deliberately small Lean file; on the artifact files they
are mutually inverse byte for byte, which is what gives the byte-identity
facts (such as Fact~\ref{fact:bytes}) their meaning.
\lean{ConwayO7.Dimacs.litOfInt, ConwayO7.Dimacs.parseDimacs,
ConwayO7.Dimacs.serializeDimacs, ConwayO7.Dimacs.parseCubes,
ConwayO7.Dimacs.serializeCubes}\leanok
\uses{def:cnf-formula, def:cnf-cube}
\end{definition}

\section{Sign-prefix cells and Kraft mass}

The cubes of the certificate are not arbitrary: they all constrain an
initial segment of one fixed sequence of variables.  Coverage therefore
reduces to a purely combinatorial statement about finite words over
$\{0,1\}$.

\begin{definition}[Sign prefixes and cells]\label{def:prefix}
Fix a \emph{branching order}: a sequence of $D$ distinct variables, in the
certificate the sequence $x_{15}, x_{16}, \dots, x_{214}$ (so $D = 200$).
To a word $s \in \{0,1\}^{\le D}$ associate the cube asserting
$x_{14+j} = s_j$ for $1 \le j \le |s|$; its satisfaction set, the
\emph{cell} of $s$, consists of all valuations extending the partial
valuation $s$ along the branching order.  A family $S$ of words is
\emph{prefix-free} if no member is a prefix of another (in particular the
members are pairwise distinct).
\lean{ConwayO7.cubeOfSigns, ConwayO7.PrefixFree}\leanok
\uses{def:cnf-cube}
\end{definition}

\begin{lemma}[Cells are prefix sets]\label{lem:cov-cell-prefix}
Write $w(a) \in \word{D}$ for the \emph{branch word} of a valuation $a$:
the values of $a$ along the branching order.  A valuation $a$ lies in the
cell of a word $s$ (with $|s| \le D$) if and only if $s$ is a prefix of
$w(a)$.  Consequently two cells are either nested or disjoint, and the
cells of a prefix-free family are pairwise disjoint.
\lean{ConwayO7.cubeOfSigns_eval_iff}\leanok
\uses{def:prefix, def:cnf-valuation}
\end{lemma}
\begin{proof}[Proof sketch]
Induction on $s$ along the branching order.
\end{proof}

\begin{definition}[Kraft weight]\label{def:cov-kraft}
The \emph{Kraft weight} at depth $D$ of a family $S$ of words of length
$\le D$ is the integer
$W_D(S) = \sum_{s \in S} 2^{\,D - |s|}$,
i.e.\ $2^D$ times the usual Kraft mass $\sum_s 2^{-|s|}$.  The cell of $s$
meets exactly $2^{\,D-|s|}$ of the $2^D$ branch words, so $W_D(S)$ is the
total number of branch words covered, counted with multiplicity.
\lean{ConwayO7.kraftWeight}\leanok
\uses{def:prefix}
\end{definition}

\begin{lemma}[Splitting at the root]\label{lem:cov-kraft-split}
For $b \in \{0,1\}$ let $S_b = \{ t : b\,t \in S \}$ be the family of tails
of the members of $S$ beginning with $b$.  If the empty word is not in $S$,
then
$W_{D+1}(S) = W_D(S_0) + W_D(S_1)$;
moreover each $S_b$ inherits prefix-freeness and the length bound from $S$.
\lean{ConwayO7.kraftWeight_split, ConwayO7.branchTails,
ConwayO7.prefixFree_branchTails}\leanok
\uses{def:cov-kraft, def:prefix}
\end{lemma}
\begin{proof}[Proof sketch]
Each $s \in S$ starts with exactly one bit $b$ and contributes
$2^{(D+1)-|s|} = 2^{D-|s_{\ge 2}|}$ to the $b$-branch.  Prefix relations
between words with the same first bit correspond exactly to prefix
relations between their tails.
\end{proof}

\begin{lemma}[Kraft's inequality]\label{lem:cov-kraft-ineq}
A prefix-free family $S \subseteq \{0,1\}^{\le D}$ satisfies
$W_D(S) \le 2^D$.
\lean{ConwayO7.kraftWeight_le}\leanok
\uses{def:cov-kraft, def:prefix}
\end{lemma}
\begin{proof}[Proof sketch]
\uses{lem:cov-kraft-split}
Induction on $D$.  If the empty word belongs to $S$ then prefix-freeness
forces $S = \{\epsilon\}$, of weight exactly $2^D$; otherwise split at the
root by Lemma~\ref{lem:cov-kraft-split} and apply the inductive bound to
each branch.
\end{proof}

\begin{lemma}[Kraft equality and maximality]\label{lem:kraft}
A prefix-free family $S \subseteq \{0,1\}^{\le D}$ with full Kraft weight
$W_D(S) = 2^D$ is a \emph{maximal antichain} of the binary tree: every
$w \in \word{D}$ has a (necessarily unique) member of $S$ as a prefix.
Hence the corresponding cells \emph{partition} the valuation space.
\lean{ConwayO7.covers_of_prefixFree_kraft}\leanok
\uses{def:prefix, def:cov-kraft}
\end{lemma}
\begin{proof}[Proof sketch]
\uses{lem:cov-kraft-split, lem:cov-kraft-ineq}
Induction on $D$, following the given word $w$ bit by bit.  If
$\epsilon \in S$ we are done.  Otherwise split at the root: by
Lemma~\ref{lem:cov-kraft-split} the two branch weights sum to
$2^{D+1} = 2^D + 2^D$, and by Kraft's inequality neither exceeds $2^D$, so
\emph{both} branches again have full weight; recurse into the branch
selected by the first bit of $w$.  Uniqueness and disjointness follow from
prefix-freeness via Lemma~\ref{lem:cov-cell-prefix}.
\end{proof}

\begin{lemma}[Coverage of the valuation space]\label{lem:cov-valuation-cover}
Let $S$ be a prefix-free family of words of length $\le D$ with full Kraft
weight $2^D$.  Then every valuation lies in the cell of some (unique)
member of $S$.
\lean{ConwayO7.cube_cover_of_signs}\leanok
\uses{def:prefix, def:cnf-valuation}
\end{lemma}
\begin{proof}[Proof sketch]
\uses{lem:kraft, lem:cov-cell-prefix}
Apply Lemma~\ref{lem:kraft} to the branch word $w(a)$ of the given
valuation $a$ and translate back through Lemma~\ref{lem:cov-cell-prefix}.
\end{proof}

\section{An executable coverage check}

The concrete family has $98{,}536$ members, so verifying prefix-freeness
literally --- all ${\sim}5 \cdot 10^9$ pairs --- is out of reach.  Instead
the family is sorted lexicographically and only \emph{adjacent} pairs are
compared; the block structure of the lexicographic order makes this
sufficient.

\begin{definition}[Lexicographic order on words]\label{def:cov-lex}
Order $\{0,1\}^{\le D}$ by $s \le t$ iff at the first position where they
differ (reading left to right) $s$ carries $0$ and $t$ carries $1$, or $s$
is a prefix of $t$.  This relation is total, transitive and antisymmetric:
a linear order on words.
\lean{ConwayO7.lexLe, ConwayO7.lexLe_total, ConwayO7.lexLe_trans,
ConwayO7.lexLe_antisymm}\leanok
\end{definition}

\begin{lemma}[Block property of the lexicographic order]\label{lem:cov-lex-block}
Extensions of a fixed word $s$ form an \emph{interval}: if $s$ is a prefix
of $t$ and $s \le u \le t$ in the lexicographic order, then $s$ is a prefix
of $u$.  (In particular $s \le t$ whenever $s$ is a prefix of $t$.)
\lean{ConwayO7.prefix_of_lexLe_of_lexLe, ConwayO7.lexLe_of_prefix}\leanok
\uses{def:cov-lex}
\end{lemma}
\begin{proof}[Proof sketch]
Simultaneous induction on the three words: matching leading bits are
stripped, and a strict inequality at the head contradicts one of the two
order hypotheses.
\end{proof}

\begin{definition}[The adjacent-pairs test]\label{def:cov-adjacent}
For a finite sequence of words, the \emph{adjacent-pairs test} checks, in
one pass, that no term is a prefix of its immediate successor.
\lean{ConwayO7.adjacentOk}\leanok
\uses{def:prefix}
\end{definition}

\begin{lemma}[Soundness of the adjacent-pairs test]\label{lem:cov-adjacent-sound}
If a sequence of words is sorted for the lexicographic order and passes the
adjacent-pairs test, then the underlying family is prefix-free.
\lean{ConwayO7.pairwise_of_sorted_adjacent}\leanok
\uses{def:cov-adjacent, def:cov-lex, def:prefix}
\end{lemma}
\begin{proof}[Proof sketch]
\uses{lem:cov-lex-block}
Suppose $s$ precedes $u$ in the sorted sequence and $s$ is a prefix of $u$,
with immediate successor $t$ of $s$.  Then $s \le t \le u$, so by the block
property (Lemma~\ref{lem:cov-lex-block}) $s$ is already a prefix of its
neighbour $t$ --- excluded by the test.  The reverse prefix relation would
force $s = u$ by antisymmetry.
\end{proof}

\begin{definition}[The coverage check]\label{def:cov-checkcover}
Given a depth $D$ and a family $S$ of words, the \emph{coverage check} is
the finite computation verifying:
(i) every member of $S$ has length $\le D$;
(ii) the family, sorted lexicographically (by merge sort), passes the
adjacent-pairs test;
(iii) the Kraft weight satisfies $W_D(S) = 2^D$ exactly, in integer
arithmetic.
\lean{ConwayO7.checkCover}\leanok
\uses{def:cov-adjacent, def:cov-lex, def:cov-kraft}
\end{definition}

\begin{lemma}[Soundness of the coverage check]\label{lem:cov-checkcover-sound}
If the coverage check succeeds on $(D, S)$, then every member of $S$ has
length $\le D$, the family $S$ is prefix-free, and $W_D(S) = 2^D$ --- the
three hypotheses of Lemma~\ref{lem:cov-valuation-cover}.
\lean{ConwayO7.cover_of_checkCover}\leanok
\uses{def:cov-checkcover, def:prefix, def:cov-kraft}
\end{lemma}
\begin{proof}[Proof sketch]
\uses{lem:cov-adjacent-sound}
The sorted sequence is a permutation of $S$ and is pairwise ordered (both
facts are standard properties of merge sort, available in the library), so
Lemma~\ref{lem:cov-adjacent-sound} yields pairwise prefix-freeness, which
transfers across the permutation; the length and weight conditions are
checked directly.
\end{proof}

\begin{theorem}[Refutation by cells]\label{thm:coverage}
Let $F$ be a CNF formula, fix a branching order of length $D$, and let
$S \subseteq \{0,1\}^{\le D}$ be prefix-free with full Kraft weight
$W_D(S) = 2^D$.  If for every $s \in S$ the formula
$F \wedge (\text{cell } s)$ is unsatisfiable, then $F$ is unsatisfiable.
\lean{ConwayO7.unsat_of_prefix_cover}\leanok
\uses{def:prefix, def:cov-kraft, def:cnf-withcube, def:cnf-unsat}
\end{theorem}
\begin{proof}[Proof sketch]
\uses{thm:cnf-unsat-of-cubes, lem:cov-valuation-cover}
By Lemma~\ref{lem:cov-valuation-cover} the cells of $S$ cover the valuation
space; conclude by the composition theorem
(Theorem~\ref{thm:cnf-unsat-of-cubes}).
\end{proof}

\begin{corollary}[Refutation by checked cells]\label{cor:cov-checked}
If the coverage check succeeds on $(D, S)$ and $F \wedge (\text{cell } s)$
is unsatisfiable for every $s \in S$, then $F$ is unsatisfiable.  The
entire coverage side of the certificate thus reduces to one finite
computation on the concrete data.
\lean{ConwayO7.unsat_of_checked_cover}\leanok
\uses{thm:coverage, def:cov-checkcover}
\end{corollary}
\begin{proof}[Proof sketch]
\uses{lem:cov-checkcover-sound}
Combine Lemma~\ref{lem:cov-checkcover-sound} with
Theorem~\ref{thm:coverage}.
\end{proof}

\section{The distributed family}

\begin{definition}[The certified tiling data]\label{def:cov-tiling}
The consolidated cell list of the artifact (one cube per certificate, in
manifest order) is embedded byte for byte into the Lean development and
parsed by the maps of Definition~\ref{def:cnf-dimacs}, yielding the list of
\emph{tiling cubes}; the \emph{branching order} is the sequence
$x_{15}, \dots, x_{214}$ of $200$ variables; the \emph{sign words} are the
words in $\{0,1\}^{\le 200}$ obtained by reading off the signs of each
cube's literals.  No processing external to Lean intervenes between the
hash-pinned file and these terms.
\lean{ConwayO7.cubesAllTxt, ConwayO7.tilingCubes, ConwayO7.tilingOrder,
ConwayO7.tilingSigns}\leanok
\uses{def:cnf-cube, def:cnf-dimacs}
\end{definition}

\begin{fact}[Faithfulness of the sign-word model]\label{fact:cov-tiling-signs}
Rebuilding each cube from its sign word along the branching order
reproduces the parsed cube list exactly: every distributed cube really is a
sign pattern on an initial segment of consecutive variables starting at
$x_{15}$.  (If the file contained any cube of another shape, this
computation would fail.)
\lean{ConwayO7.tilingSigns_eq}\leanok
\uses{def:cov-tiling, def:prefix}
\end{fact}

\begin{fact}[Cardinality]\label{fact:cov-tiling-count}
The distributed family consists of exactly $98{,}536$ sign words, of
lengths between $12$ and $200$.
\lean{ConwayO7.tiling_count}\leanok
\uses{def:cov-tiling}
\end{fact}

\begin{fact}[The distributed family tiles the valuation space]\label{fact:tiling}
The coverage check of Definition~\ref{def:cov-checkcover} succeeds, at
depth $200$, on the $98{,}536$ sign words: the family is prefix-free with
full Kraft weight $2^{200}$, so by Lemma~\ref{lem:kraft} its cells
partition the valuation space --- a perfect tiling.  This is established by
executing the verified check, inside Lean's compiled evaluator, on the
embedded data.
\lean{ConwayO7.o7Tiling_checkCover}\leanok
\uses{def:cov-tiling, def:cov-checkcover, def:prefix}
\end{fact}

%=====================================================================
